\chapter{Antecedentes y justificación}

\section{Antecedentes de la carrera en la UCR}

La formación universitaria en matemática en Costa Rica tiene sus raíces en la Escuela de Ciencias de la Universidad de Costa Rica, creada en 1940, donde la disciplina atendía principalmente las necesidades de otras carreras y la formación de profesores para enseñanza media. En 1957, con la Reforma Universitaria de Rodrigo Facio, se integró el Departamento de Física y Matemáticas, antecedente directo de la carrera profesional en Matemática. En 1971 se creó el Departamento de Matemáticas y, producto del Tercer Congreso Universitario (1972--1973), en 1974 el departamento se convirtió en Escuela de Matemática, con autonomía para impartir y administrar su propia carrera.

Desde entonces, la Escuela ha ofrecido Bachillerato y Licenciatura en Matemática en la Sede Rodrigo Facio (Campus Universitario). Entre los principales hitos curriculares e institucionales se destacan: la Licenciatura en Matemática (1967--1972), el plan con énfasis purista vigente entre 1973 y 1991 (124 créditos de bachillerato y 30 de licenciatura), las reformas de 1992 (cursos optativos interdisciplinarios y ajustes al primer año), el Programa de Maestría en Matemáticas (desde 1980), la creación de los centros CIMM y CIMPA en la década de 1990, y la consolidación de la investigación y la divulgación mediante publicaciones como la \emph{Revista de Matemática: Teoría y Aplicaciones}.

En 2017 se abrió la carrera de Educación Matemática. Un año después inició, de manera conjunta con Ciencias Actuariales, un proceso de reestructuración curricular que en 2019 se separó por carrera. Entre 2018 y 2020, una subcomisión de la Comisión de Docencia de la Escuela---con apoyo del Centro de Enseñanza y Aprendizaje (CEA) y mandato de la Vicerrectoría de Docencia---elaboró el perfil de salida y los marcos referenciales que fundamentan la presente propuesta. Ese trabajo incorporó entrevistas, grupos focales, análisis estructurales de la actividad, talleres con docentes y estudiantes, y consultas a la comunidad académica.

Actualmente la Escuela de Matemática cuenta con 128 profesores (41 con grado de doctor), organizados en tres departamentos: Matemática Pura y Ciencias Actuariales, Educación Matemática, y Matemática Aplicada. El Bachillerato y Licenciatura en Matemática se imparten desde el Departamento de Matemática Pura y Ciencias Actuariales. La propuesta curricular que se presenta responde a la necesidad de actualizar un plan cuyo último rediseño integral data de hace más de tres décadas, en un contexto marcado por la digitalización, la interdisciplinariedad y nuevas demandas del mercado laboral y de la investigación.


\section{Justificación socioprofesional}

La Matemática ocupa un lugar central en el desarrollo científico, tecnológico y social contemporáneo. La digitalización de la información, el avance de la inteligencia artificial, el manejo de datos a gran escala y la modelación de fenómenos complejos han ampliado tanto el objeto de estudio de la disciplina como los ámbitos en que el matemático puede incidir. La UNESCO proclamó el 14 de marzo como Día Internacional de las Matemáticas (2019) y subraya su papel en el logro de los Objetivos de Desarrollo Sostenible, en la formación de capacidades para enfrentar desafíos como el cambio climático, la salud pública y la innovación tecnológica.

En este marco, la carrera de Matemática de la UCR debe formar profesionales competentes para la academia, el sector productivo y la interacción con otras disciplinas. La propuesta de plan de estudios articula cinco áreas curriculares---habilidades operacionales, pensamiento abstracto, formalismo, pensamiento algorítmico y habilidades investigativas---y nuevos espacios formativos (modelación matemática, seminario, comunicación en las ciencias, énfasis en matemática pura y aplicada) que responden a las demandas actuales y futuras del ejercicio profesional.


\subsection{Caracterización profesional}

\subsubsection{Radiografía laboral}

A nivel internacional, las ocupaciones vinculadas con la matemática muestran una demanda creciente. La Oficina de Trabajo Estadístico de Estados Unidos proyectó un aumento del 28\,\% en empleos matemáticos entre 2016 y 2026, superior al promedio de otras ocupaciones, impulsado por el \emph{big data} y la analítica. En rankings como los de CareerCast (2016), figuras como analista de datos, estadístico, matemático y actuario ocupan posiciones destacadas entre las mejores profesiones; carreras con alto contenido matemático (ciencias de la computación, ingeniería, economía, meteorología, entre otras) también se ubican en los primeros lugares.

En Costa Rica, el matemático encuentra oportunidades en instituciones de educación superior e investigación, centros como el CIMM y el CIMPA, organismos públicos y privados que requieren modelación, análisis cuantitativo y apoyo a la toma de decisiones, así como en sectores como finanzas, banca, economía, tecnología y servicios basados en datos. La separación de funciones actuariales hacia la carrera de Ciencias Actuariales y la especialización docente en Educación Matemática redefinen el espacio ocupacional del matemático, orientándolo hacia la investigación, la docencia universitaria, la modelación, el análisis de datos y el trabajo interdisciplinario.

La formación en Matemática desarrolla capacidades analíticas, rigor argumentativo y resolución de problemas que son valoradas transversalmente por empleadores, conforme a evidencia en neurociencia cognitiva y a la propia naturaleza de la disciplina como lenguaje universal de las ciencias.


\subsubsection{Situación de la población graduada}

En el marco del proceso de reestructuración 2018--2020, la Escuela de Matemática recopiló información sobre las prácticas profesionales vigentes mediante entrevistas a especialistas, análisis estructurales de la actividad por identidad matemática, talleres con docentes y estudiantes avanzados, y consultas a la comunidad académica. Esos insumos permitieron caracterizar los escenarios de inserción y las funciones que desempeñan los matemáticos formados en la UCR.

De manera consistente con ese análisis, la población graduada se inserta principalmente en la docencia e investigación universitaria, en la modelación matemática de fenómenos de otras disciplinas y, de manera creciente, en el sector productivo (finanzas, banca, análisis de datos, consultoría cuantitativa). La dificultad para acceder a empleos directamente vinculados con la carrera varía según el ámbito: persiste demanda en academia e investigación, mientras que la enseñanza en secundaria y las labores actuariales han quedado en mayor medida a cargo de profesionales de Educación Matemática y Ciencias Actuariales, respectivamente.

La valoración de la formación recibida---especialmente respecto del perfil de egreso---constituyó insumo central para la elaboración del perfil de salida propuesto, que por primera vez explicita las competencias del matemático del siglo XXI. En la etapa de formalización curricular, la unidad académica complementará este apartado con el seguimiento institucional a las últimas cohortes graduadas, conforme a los lineamientos del Observatorio Laboral de Profesiones (OLAP) y otras fuentes nacionales.


\subsubsection{Campos de inserción laboral}

Los principales campos de inserción identificados para el matemático formado en la UCR son:

\begin{itemize}
  \item \textbf{Investigación académica} (teórica y aplicada), en universidades, centros de investigación y redes científicas nacionales e internacionales.
  \item \textbf{Docencia universitaria}, en programas de pregrado y posgrado en matemática y disciplinas afines.
  \item \textbf{Modelación matemática}, en ingeniería, física, biología, epidemiología, ciencias ambientales y otras áreas que requieren traducir problemas reales al lenguaje matemático.
  \item \textbf{Análisis de datos y ciencia de datos}, en empresas, instituciones financieras, organismos públicos y proyectos de investigación interdisciplinaria.
  \item \textbf{Sector productivo e industrial}, en roles de consultoría cuantitativa, optimización, criptografía, simulación y apoyo a la innovación tecnológica.
  \item \textbf{Trabajo interdisciplinario}, como puente entre equipos de distintas disciplinas en proyectos científicos, tecnológicos o de política pública.
\end{itemize}


\subsubsection{Funciones o tareas profesionales}

Para cada campo de inserción, el matemático desempeña funciones que, sintetizadas a partir del análisis estructural de la actividad realizado en la Escuela, incluyen:

\begin{itemize}
  \item Estudiar, clasificar y relacionar estructuras matemáticas (algebraicas, geométricas, analíticas, topológicas, probabilísticas, etc.) para comprender y resolver problemas.
  \item Formular conjeturas, demostrar resultados y contribuir a la creación de conocimiento matemático original.
  \item Traducir problemas de la vida real o de otras disciplinas al lenguaje matemático, seleccionar herramientas adecuadas y comunicar resultados en lenguaje interdisciplinario.
  \item Modelar fenómenos naturales, sociales o tecnológicos mediante ecuaciones, simulaciones y métodos numéricos o estocásticos.
  \item Analizar e interpretar datos, identificar patrones y apoyar la toma de decisiones basada en evidencia.
  \item Enseñar matemática en el ámbito universitario, diseñar experiencias de aprendizaje y orientar procesos de formación e investigación.
  \item Divulgar la producción científica mediante artículos, informes, presentaciones y participación en congresos y seminarios.
\end{itemize}

Estas funciones orientan la definición del perfil de egreso y la selección de contenidos, competencias y tipos de curso de la malla propuesta.


\subsection{Contribución al desarrollo del país}

La Matemática contribuye al proyecto humanista de la Universidad de Costa Rica al formar profesionales capaces de producir y aplicar conocimiento riguroso al servicio de la sociedad. Las matemáticas son herramienta para el desarrollo sostenible: modelan cambios ambientales y epidemiológicos, optimizan recursos en salud pública y transporte, sustentan la criptografía y las comunicaciones seguras, y apoyan la inteligencia artificial y el aprendizaje automático con aplicaciones en medicina, agricultura, energía y gobernanza.

En el plano nacional, la Escuela de Matemática ha sido pionera en la introducción de las matemáticas modernas en la enseñanza secundaria (Conferencia Interamericana de Educación Matemática, Bogotá, 1961) y mantiene un rol activo en la producción de conocimiento mediante el CIMM, el CIMPA y la Maestría en Matemáticas. La investigación y divulgación matemática enriquecen la cultura científica del país, fortalecen la formación de cuadros docentes e investigadores y generan capacidades para la innovación en sectores estratégicos de la economía costarricense.

La propuesta curricular responde además a compromisos institucionales y a marcos internacionales como los Objetivos de Desarrollo Sostenible y las prioridades de la UNESCO en ciencias básicas, educación matemática e igualdad de género. Al incorporar énfasis en matemática pura y aplicada, modelación, computación y comunicación científica, el plan busca que los graduados contribuyan a la calidad de vida de las personas, al crecimiento del capital humano avanzado y a la generación de oportunidades laborales en ámbitos emergentes de alta demanda social.


\subsection{Prospectiva de la profesión}

El análisis de prácticas profesionales realizado en la Escuela distingue tendencias emergentes, dominantes y decadentes que orientan la proyección de la profesión en la próxima década.

Entre las \textbf{prácticas emergentes} se destacan: la investigación fuera de la academia en disciplinas afines (finanzas, banca, economía, tecnología); el análisis y la ciencia de datos a gran escala; y la experimentación científica apoyada en simulación y recursos computacionales. Entre las \textbf{prácticas dominantes} permanecen la investigación académica, la docencia universitaria, la modelación de fenómenos interdisciplinarios y el trabajo colaborativo con otras ciencias. Entre las \textbf{prácticas decadentes}, en el contexto costarricense, figuran las labores actuariales propias del matemático (asumidas por actuarios) y la docencia en enseñanza secundaria sin formación en Educación Matemática.

A nivel internacional, las mallas de referencia incorporan mayor énfasis en computación, modelación, probabilidad aplicada, comunicación matemática y proyectos de investigación formativa. La propuesta de la UCR alinea la formación con esas tendencias mediante cursos de matemática computacional, modelación matemática, seminario, comunicación en las ciencias y rutas de énfasis en matemática pura y aplicada, sin abandonar la solidez en álgebra, análisis, geometría y topología que caracteriza la tradición formativa de la Escuela.

En los próximos diez años se prevé que el matemático costarricense deba combinar rigor teórico con versatilidad aplicada, capacidad de trabajo en equipos interdisciplinarios, dominio de herramientas computacionales y habilidades de comunicación científica. La presente reestructuración curricular busca anticipar esos hitos y preparar graduados capaces de incursionar en instituciones de investigación de prestigio internacional o en el sector productivo e industrial, en coherencia con la declaración de propósitos de la carrera.
